\section{Metode og evalueringsgrunnlag}
Prosjektet er gjennomført som et design- og implementasjonsprosjekt med
empirisk evaluering. Først ble \gls{hdo}s kravspesifikasjon analysert for å
identifisere funksjonelle og ikke-funksjonelle krav. Deretter ble løsningen
designet, implementert og evaluert i flere iterasjoner mot \gls{hdo}s testmiljø.

Evalueringen som besvarer forskningsspørsmålet bygger på fire typer grunnlag.
\emph{Krav-evaluering} sammenstiller hvert krav i kravspesifikasjonen mot
hvor det er implementert og hvordan det er testet, og besvarer i hovedsak
dimensjon~1 og dimensjon~3. \emph{Arkitektur-evaluering} vurderer hvor
backend-spesifikk logikk er plassert og hvilke endringer som faktisk kreves
for å støtte en ny videomotor; dette ble undersøkt empirisk ved å legge til
to ekstra providere etter at hovedimplementasjonen var ferdig (se
seksjon~\ref{sec:provider-validering}), og besvarer i hovedsak
dimensjon~2. \emph{Funksjonell verifisering} ble gjort gjennom 133
automatiserte enhets- og integrasjonstester, samt et ende-til-ende-skript
som kjører mot \gls{hdo}s faktiske AntMedia- og \gls{nhn}-testmiljøer.
\emph{Drifts-evaluering} dekker containerisert kjøring, helsesjekker,
Prometheus-metrikker og Grafana-dashboard, og bidrar primært til
dimensjon~3.

I tillegg ble det innhentet en strukturert sluttevaluering fra \gls{hdo} etter
overlevering, der to representanter fra oppdragsgiver vurderte leveransen
mot kravområdene på en skala fra 1 til 5. Sluttevalueringen er gjengitt
i sin helhet i vedlegg~\ref{app:hdo-evaluering} og brukes som ekstern
validering på tvers av alle tre dimensjoner. Datagrunnlaget er for
begrenset til statistiske slutninger, men gir et kvalitativt bilde av
hvor godt leveransen møter behovet hos oppdragsgiver.